% =====================================================================
% الفصل الأول: الإطار المفاهيمي: الشغل المؤقت والمناولة والتمييز بينهما
% =====================================================================

\chapter{الإطار المفاهيمي: الشغل المؤقت والمناولة والتمييز بينهما}
\label{chap:chapitre1}

\section*{مقدمة الفصل}

يعد الضبط المفاهيمي منطلقاً أساسياً وحجر الزاوية في أي دراسة قانونية واجتماعية معمقة؛ إذ إن الخلط بين المفاهيم في حقل العلاقات المهنية لا يقف عند حدود التوصيف النظري، بل ينعكس مباشرة على تحديد القواعد الآمرة الواجبة التطبيق، وتعيين أطراف المسؤولية، ورسم معالم الحماية المقررة للطرف الضعيف في الرابطة العقدية وهو الأجير[cite: 1, 2]. 

ويبرز هذا التداخل جلياً في الممارسة العملية بين تقنيتي \textbf{الشغل المؤقت} (\textit{Travail temporaire}) و\textbf{المناولة أو المقاولة من الباطن} (\textit{Sous-traitance})؛ فهما يشتركان ظاهرياً في إقحام طرف ثالث وسيط بين العامل والمنشأة المستفيدة من جهده العضلي أو الفكري[cite: 2]. غير أن هذا التشابه الشكلي يخفي تبايناً جوهرياً في الفلسفة القانونية ومحل الالتزام وطبيعة التبعية[cite: 1, 2].

يهدف هذا الفصل إلى تأصيل مفهوم الشغل المؤقت وتفكيك هندسته التعاقدية الثلاثية (المبحث الأول)، ثم تحليل مفهوم المقاولة من الباطن وشروط مشروعيتها (المبحث الثاني)، قبل استعراض معايير التمييز النسقية بينهما وبين عقد الشغل المحدد المدة (CDD)، وتنزيل ذلك تحليلياً على واقع النازلة الميدانية المعروضة في سياق هذه الدراسة (المبحث الثالث)[cite: 1, 2, 3].

% ---------------------------------------------------------------------
\section{مفهوم الشغل المؤقت وهندسته التعاقدية}

\subsection{التعريف القانوني لمقاولة التشغيل المؤقت}
عرف المشرع المغربي مقاولة التشغيل المؤقت في المادة 495 من مدونة الشغل بأنها: \textit{«كل شخص اعتباري، مستقل عن السلطة العمومية، يقتصر عمله على ممارسة النشاط المشار إليه في الفقرة (ج) من المادة 477، والمتمثل في تشغيل أجراء بهدف وضعهم، مؤقتاً، رهن إشارة شخص ثالث يسمى "المستعمل"، يحدد مهامهم ويراقب تنفيذها»}[cite: 2].

وانطلاقاً من هذا النص الصريح، يتأسس مفهوم الشغل المؤقت على أركان قانونية آمرة تهدف إلى ضبط هذه المهنة ومنع انحرافها:
\begin{enumerate}
    \item \textbf{الشخصية الاعتبارية المسبقة:} حصر المشرع مزاولة هذا النشاط في الأشخاص الاعتباريين (شركات تجارية)[cite: 2, 3]، مانعاً الأشخاص الذاتيين (الطبيعيين) من ممارسته[cite: 3]. والغاية من هذا الشرط هي ضمان ملاءة مالية ومسك محاسبة منتظمة تتيح مراقبة التدفقات المالية وقدرة الشركة على الوفاء بأجور الأجراء واشتراكاتهم الاجتماعية[cite: 2, 3].
    \item \textbf{مبدأ حصرية النشاط (\textit{Principe d'exclusivité}):} أوجبت المادة 495 أن «يقتصر عملها» على نشاط وضع الأجراء[cite: 2]. فلا يجوز قانوناً لشركة التشغيل المؤقت أن تجمع بين هذا النشاط وأنشطة تجارية أخرى، كالمقاولات العامة أو الخدمات أو التوريدات؛ وذلك لقطع الطريق أمام استغلال الأجراء في أنشطة موازية أو خلط الذمم المالية والمصالح التجارية[cite: 2].
    \item \textbf{الطابع الاستثنائي والعارض للعمل («المهمة»):} لا يمارس الأجير المؤقت عملاً دائماً ومستمراً، بل يعين لإنجاز «مهمة» (\textit{Mission}) محددة الأجل والهدف، ومحصورة في الحالات التي عددها المشرع استثناءً في المادة 496 من المدونة[cite: 2, 3].
\end{enumerate}

\subsection{العلاقة الثلاثية وتجزئة التبعية القانونية}
تتميز علاقة الشغل المؤقت بتركيبتها الثلاثية غير النمطية، إذ تنفصم فيها الثنائية الكلاسيكية لعقد الشغل (أجير - مشغل) لتتحول إلى شبكة تعاقدية تضم ثلاثة فاعلين أساسيين تحكمهم رابطتان عقديتان[cite: 2]:

\begin{itemize}
    \item \textbf{الرابطة العقدية الأولى (عقد الشغل المؤقت):} وهو عقد يربط بين مقاولة التشغيل المؤقت والأجير[cite: 2]. وبموجبه تعد المقاولة هي \textbf{المشغل القانوني الحصري}؛ فهي التي تتحمل عبء أداء الأجر، والتعويض عن العطل المؤدى عنها، والتصريح الإجباري بالأجير لدى الصندوق الوطني للضمان الاجتماعي، وممارسة السلطة التأديبية والجزائية كالفصل أو الإنذار[cite: 2, 3].
    \item \textbf{الرابطة العقدية الثانية (عقد وضع الأجير رهن الإشارة):} وهو عقد تجاري مدني يربط مقاولة التشغيل المؤقت بالمقاولة المستعملة[cite: 2, 3]. وقد اشترطت المادة 499 كتابته وبيان سبب اللجوء، ومدة المهمة، ومكانها، والمقابل المالي المتفق عليه لتوفير اليد العاملة[cite: 2, 3].
\end{itemize}

وينجم عن هذه الهندسة ما يسميه الفقه الاجتماعي بـ \textbf{«تجزئة التبعية القانونية»} (\textit{Démembrement du lien de subordination})؛ حيث تنشطر عناصر التبعية إلى شقين:
\begin{itemize}
    \item \textbf{تبعية قانونية وإدارية:} تظل لصيقة بمقاولة التشغيل المؤقت (إبرام العقد، الأجر، التأديب، إنهاء العلاقة التعاقدية)[cite: 2].
    \item \textbf{تبعية تقنية ووظيفية وميدانية:} تنتقل فعلياً إلى المقاولة المستعملة؛ إذ يمارس مسؤولوها سلطة الإشراف والمراقبة اليومية، وتوجيه التعليمات التقنية، وتحديد وتائر العمل وأوقاته، وإلزام الأجير بقواعد النظام الداخلي وتدابير السلامة والوقاية المعمول بها داخل المنشأة[cite: 2, 3].
\end{itemize}

% ---------------------------------------------------------------------
\section{مفهوم المناولة (المقاولة من الباطن)}

\subsection{التأصيل القانوني لعقد المقاولة من الباطن}
نصت المادة 86 من مدونة الشغل على أن المقاولة من الباطن هي: \textit{«عقد مكتوب يكلف بمقتضاه مقاول أصلي مقاولاً من الباطن بأن يقوم له بشغل من الأشغال، أو بأن ينجز له خدمة من الخدمات، كلما كان ذلك في صالح المقاولة الأصلية ودون إخلال بحقوق أجرائها»}[cite: 2, 3].

والمناولة في جوهرها الاقتصادي والقانوني هي أسلوب من أساليب التنظيم الصناعي والإنتاجي تتوخى من خلاله المقاولة الأصلية تفويض أجزاء من نشاطها إلى مقاولين مستقلين يمتلكون المعرفة التقنية والوسائل اللوجستية التي تفتقر إليها المنشأة الأصلية، وتتخذ في التطبيق شكلين:
\begin{itemize}
    \item \textbf{مناولة التخصص (\textit{Sous-traitance de spécialité}):} حيث تفتقر المقاولة الأصلية للخبرة التقنية أو التجهيزات الدقيقة، فتعمد إلى إسناد أشغال الصيانة المعقدة، أو الأنظمة المعلوماتية، أو الأمن والحراسة والنظافة إلى شركات مختصة تمتلك الدراية والمهارات اللازمة.
    \item \textbf{مناولة السعة (\textit{Sous-traitance de capacité}):} وتلجأ إليها المنشأة عندما تعجز وسائلها الذاتية عن استيعاب حجم الطلبات المتزايدة مؤقتاً، فتسند جزءاً من التصنيع لمقاول خارجي ينجزه في ورشاته الخاصة أو داخل منشأة المقاول الأصلي بوسائل إشراف مستقلة.
\end{itemize}

\subsection{أركان العقد وشروط مشروعيته}
لكي يظل عقد المناولة صحيحاً ومشروعاً ولا ينحرف إلى عملية اتجار غير مشروع بقوة العمل، يجب أن تتوافر فيه جملة من الشروط والأركان الجوهرية[cite: 2]:
\begin{enumerate}
    \item \textbf{شكلية الكتابة:} فرضت المادة 86 إفراغ العقد في محرر مكتوب لبيان حدود العملية وطبيعة المهام[cite: 2, 3].
    \item \textbf{محل العقد (تحقيق نتيجة متخصصة):} يلتزم المقاول من الباطن بإنجاز شغل أو أداء خدمة محددة المعالم تقنياً تنتهي بتحقيق نتيجة ملموسة ومستقلة، وليس مجرد تقديم وتوريد أجراء أو وضع ساعات عمل رهن تصرف المقاول الأصلي[cite: 2].
    \item \textbf{الاستقلال الإداري والتقني الكامل:} يظل المقاول من الباطن هو المشغل الحصري لأجرائه قانونياً وفعلياً؛ فهو الذي يمارس عليهم سلطة الإدارة والتوجيه والإشراف والتأديب عبر أطره ورؤساء ورشاته، دون أن يتدخل المقاول الأصلي في توجيه عمال المناولة بصورة مباشرة[cite: 2].
    \item \textbf{استخدام وسائل ومعدات ذاتية:} يفترض في المقاول من الباطن أن يقدم أدوات العمل ومعداته ومواده الأولية، ما لم تتطلب طبيعة الخدمة استعمال آلات المنشأة الأصلية بمقتضى بند تعاقدي صريح ومبرر تقنياً.
    \item \textbf{شرط الحفاظ على حقوق الأجراء الأصليين:} علق المشرع مشروعية اللجوء إلى المناولة بشرط \textit{«ألا يكون مخلاً بحقوق أجراء المقاولة الأصلية»}[cite: 2, 3]، مانعاً استعمالها لتجاوز مكتسبات الأجراء، أو كوسيلة للإغلاق الجزئي للمنشأة لتسريح العمال وتعويضهم بمقاولات باطن أقل تكلفة[cite: 2].
\end{enumerate}

% ---------------------------------------------------------------------
\section{التمييز بين الشغل المؤقت والمناولة وعقد الشغل المحدد المدة}

\subsection{المقارنة المنهجية بين الشغل المؤقت والمناولة}
لتوضيح الحدود الفاصلة بدقة بين المؤسستين، يقدم الجدول المقارن الآتي تحليلاً تفصيلياً لعناصر التمايز بينهما:

\begin{table}[h]
\centering
\small
\renewcommand{\arraystretch}{1.55}
\begin{tabularx}{\textwidth}{|p{2.8cm}|X|X|}
\hline
\textbf{وجه المقارنة} & \textbf{الشغل المؤقت (\textit{Intérim})} & \textbf{المناولة (\textit{Sous-traitance})} \\
\hline
\textbf{الأساس التشريعي} & المواد 495 وما يليها من مدونة الشغل[cite: 2, 3] & المواد 86 وما يليها من مدونة الشغل[cite: 2, 3] \\
\hline
\textbf{محل الالتزام} & وضع قوة عمل بشرية رهن إشارة المستعمل[cite: 2] & أداء شغل فني أو تقديم خدمة متخصصة متكاملة[cite: 2] \\
\hline
\textbf{طبيعة التزام المقاول} & التزام ببذل عناية (توفير عمال مؤهلين)[cite: 2] & التزام بتحقيق نتيجة محددة ومستقلة[cite: 2] \\
\hline
\textbf{السلطة والإشراف الميداني} & سلطة التوجيه والمراقبة للمقاولة المستعملة[cite: 2] & السلطة الفنية والتوجيهية حصرية لمقاول الباطن[cite: 2] \\
\hline
\textbf{وسائل العمل والتجهيزات} & تعود للمقاولة المستعملة وتوضع تحت تصرف الأجير & يوفرها المقاول من الباطن كأصل عام[cite: 2] \\
\hline
\textbf{المقابل المالي والفوترة} & فاتورة تحتسب على أساس عدد ساعات العمل والأجور & فاتورة إجمالية أو جزافية تحدد بحسب إنجاز الخدمة \\
\hline
\textbf{مدة التعاقد} & محددة بمدد قصوى حصرية بموجب المادة 500[cite: 2, 3] & ترتبط بالمدة الضرورية لإنجاز الصفقة أو المشروع[cite: 2] \\
\hline
\textbf{المسؤولية القانونية} & مسؤولية مباشرة وتضامنية في السلامة والتأمين (المادة 504)[cite: 3] & مسؤولية احتياطية مشروطة بحالات الإعسار (المادتان 89 و90)[cite: 2, 3] \\
\hline
\end{tabularx}
\caption{المقارنة النسقية والمفاهيمية بين الشغل المؤقت وعقد المناولة}
\label{tab:comparatif-detaille}
\end{table}

\subsection{التمييز عن عقد الشغل المحدد المدة (CDD)}
يمثل عقد الشغل المحدد المدة، المنظم بمقتضى المواد 16 وما يليها من مدونة الشغل، الإطار القانوني التقليدي للمرونة المباشرة[cite: 2, 3]:
\begin{itemize}
    \item \textbf{رابطة ثنائية مباشرة:} يبرم العقد مباشرة بين المشغل المستفيد والأجير، وتجتمع في المشغل كافة عناصر التبعية القانونية والفنية دون أي وساطة[cite: 2].
    \item \textbf{الاستثناء من الأصل:} يظل العقد غير محدد المدة (CDI) هو الأصل العام والمهيمن في مدونة الشغل، ولا يجوز إبرام عقد محدد المدة إلا في الحالات المحددة حصراً بنص المادة 16 والمرسوم رقم 2.19.793[cite: 3] (إحلال أجير، تزايد نشاط مؤقت، أو قطاعات موسمية محددة).
    \item \textbf{الجزاء المترتب عن التجاوز:} إذا انتهت المدة المحددة واستمر الأجير في أداء عمله دون اعتراض المشغل، يتحول العقد بقوة القانون وتلقائياً إلى عقد غير محدد المدة، وهو جزاء حمائي صارم يفتقر إليه الشغل المؤقت في بعض صوره التطبيقية بالمغرب.
\end{itemize}

\subsection{إسقاط التمييز المفاهيمي على نازلة الدراسة}
بإنزال هذه المعايير الفقهية والضوابط القانونية على واقعة النازلة التي تشكل منطلق بحثنا — والمتمثلة في تعاقد شركة صناعية كبرى مع شركة مناولة لتوفير عمال يشتغلون داخل خطوط إنتاج المصنع لسنوات تحت إشراف أطر الشركة الأصلية[cite: 1, 2] — تبرز جملة من الخلاصات التحليلية الحاسمة:

\begin{enumerate}
    \item \textbf{انتفاء الصفة الحقيقية لعقد المناولة:} إن قصر موضوع التعاقد على «توفير اليد العاملة والعمالة اليدوية» يجرد العقد من جوهر المناولة القائم على التخصص والنتيجة المستقلة[cite: 2]؛ فالشركة الأصلية لم تسند خدمة متكاملة (كنظافة المنشأة أو حراستها أو صيانة آلاتها)، بل استعارت عمالاً لدمجهم في صلب الدورة الإنتاجية للمصنع.
    \item \textbf{ثبوت ممارسة التبعية المباشرة:} إن خضوع العمال للأوامر اليومية والتوجيه الفني الصادر عن أطر ومهندسي الشركة الأصلية ينفي الاستقلال الإداري والفني لشركة المناولة[cite: 2]، مما يجعل هذه الأخيرة مجرد «مشغل صوري» أو وسيط مالي لتوزيع الأجور.
    \item \textbf{مخالفة قواعد الشغل المؤقت والمدد القانونية:} حتى لو حاول الأطراف تكييف العقد على أنه وضع يد عاملة، فإن المشرع يحظر هذا النشاط على غير مقاولات التشغيل المؤقت المرخصة[cite: 2, 3]، ويمنع استخدام الأجراء لسنوات متواصلة، إذ تحصر المادة 500 المدد القصوى للمهام في مدد لا تتجاوز 3 أو 6 أشهر[cite: 3].
    \item \textbf{إعمال مبدأ أولوية الواقع على المظهر (\textit{Primauté de la réalité}):} بموجب هذا المبدأ المستقر في الاجتهاد القضائي المقارن والمغربي والمكرس في التوصية رقم 198 لمنظمة العمل الدولية، فإن القاضي الاجتماعي لا يقف عند الأوصاف التي أطلقها الأطراف على عقودهم، بل يستنبط الحقيقة من الممارسة اليومية. وبما أن العمال يلبون حاجة قارة ودائمة، فإن العلاقة التعاقدية الحقيقية هي عقد شغل مباشر وغير محدد المدة (CDI) يربط العمال بالشركة الصناعية الأصلية[cite: 1, 2].
\end{enumerate}

\section*{خلاصة الفصل}
يخلص هذا الفصل إلى أن المشرع المغربي، وإن فتح باب المرونة من خلال الشغل المؤقت والمناولة، فإنه أحاطهما بضوابط موضوعية وشكلية دقيقة لمنع استغلالهما كأدوات للالتفاف على استقرار الشغل[cite: 1, 2]. وتكشف دراسة الواقعة الميدانية أن استمرار العمال لسنوات داخل المصنع ينفي عن العقد صفة المناولة ويحوله إلى وساطة مشبوهة تفتقد للسند القانوني. ولتبيان النصوص المنظمة والواجبة التطبيق بتفصيل، ننتقل في الفصل الثاني لدراسة الإطار القانوني للشغل المؤقت والمناولة في مدونة الشغل المغربية[cite: 2, 3].